\begin{table}[!htbp]
\centering
\small
\caption{The zero $x_{n,\alpha}$ in the second experiment ($\alpha=-1/2$, $j=2$, $\lambda=1$; Charlier $\mu=5$, $a_1=1$; Meixner $\gamma=2$, $\mu=1/2$, $a_1=2$), with $e_n$ as in Table~\ref{tab:exceptional}. Here $-1<\alpha<0$ and the zero lies to the right of $\alpha$.}
\label{tab:half}
\renewcommand{\arraystretch}{1.15}
\begin{tabular}{rcccccc}
\toprule
& \multicolumn{3}{c}{Charlier} & \multicolumn{3}{c}{Meixner}\\
\cmidrule(lr){2-4}\cmidrule(lr){5-7}
$n$ & $x_{n,\alpha}$ & $n(x_{n,\alpha}-\alpha)$ & $n^{2}e_n$ & $x_{n,\alpha}$ & $n(x_{n,\alpha}-\alpha)$ & $n^{2}e_n$\\
\midrule
$50$ & $-0.47651019$ & $1.1745$ & $8.72$ & $-0.45219010$ & $2.3905$ & $19.52$\\
$100$ & $-0.48922508$ & $1.0775$ & $7.75$ & $-0.47831271$ & $2.1687$ & $16.87$\\
$200$ & $-0.49481620$ & $1.0368$ & $7.35$ & $-0.48960334$ & $2.0793$ & $15.87$\\
$400$ & $-0.49745518$ & $1.0179$ & $7.17$ & $-0.49490364$ & $2.0385$ & $15.42$\\
$800$ & $-0.49873893$ & $1.0089$ & $7.08$ & $-0.49747624$ & $2.0190$ & $15.21$\\
$1600$ & $-0.49937225$ & $1.0044$ & $7.04$ & $-0.49874410$ & $2.0094$ & $15.10$\\
\bottomrule
\end{tabular}
\end{table}
